\documentclass[10pt,a4paper]{amsart}
\usepackage{amsmath,amssymb,amsthm,graphicx}
\usepackage{verbatim}
\newenvironment{reference}{\comment}{\endcomment}
%\newenvironment{reference}{}{}
\newenvironment{slogan}{\comment}{\endcomment}
\newenvironment{history}{\comment}{\endcomment}

% For commutative diagrams we use Xy-pic
\usepackage[all]{xy}

% We use 2cell for 2-commutative diagrams.
\xyoption{2cell}
% Theorem environments.
%
\theoremstyle{plain}
\newtheorem{theorem}[subsection]{Theorem}
\newtheorem{proposition}[subsection]{Proposition}
\newtheorem{lemma}[subsection]{Lemma}

\theoremstyle{definition}
\newtheorem{definition}[subsection]{Definition}
\newtheorem{example}[subsection]{Example}
\newtheorem{exercise}[subsection]{Exercise}
\newtheorem{situation}[subsection]{Situation}

\theoremstyle{remark}
\newtheorem{remark}[subsection]{Remark}
\newtheorem{remarks}[subsection]{Remarks}

\numberwithin{equation}{subsection}

% Macros
%
\def\lim{\mathop{\mathrm{lim}}\nolimits}
\def\colim{\mathop{\mathrm{colim}}\nolimits}
\def\Spec{\mathop{\mathrm{Spec}}}
\def\Hom{\mathop{\mathrm{Hom}}\nolimits}
\def\Ext{\mathop{\mathrm{Ext}}\nolimits}
\def\SheafHom{\mathop{\mathcal{H}\!\mathit{om}}\nolimits}
\def\SheafExt{\mathop{\mathcal{E}\!\mathit{xt}}\nolimits}
\def\Sch{\mathit{Sch}}
\def\Mor{\mathop{\mathrm{Mor}}\nolimits}
\def\Ob{\mathop{\mathrm{Ob}}\nolimits}
\def\Sh{\mathop{\mathit{Sh}}\nolimits}
\def\NL{\mathop{N\!L}\nolimits}
\def\CH{\mathop{\mathrm{CH}}\nolimits}
\def\proetale{{pro\text{-}\acute{e}tale}}
\def\etale{{\acute{e}tale}}
\def\QCoh{\mathit{QCoh}}
\def\Ker{\mathop{\mathrm{Ker}}}
\def\Im{\mathop{\mathrm{Im}}}
\def\Coker{\mathop{\mathrm{Coker}}}
\def\Coim{\mathop{\mathrm{Coim}}}

% Boxtimes
%
\DeclareMathSymbol{\boxtimes}{\mathbin}{AMSa}{"02}

%
% Macros for moduli stacks/spaces
%
\def\QCohstack{\mathcal{QC}\!\mathit{oh}}
\def\Cohstack{\mathcal{C}\!\mathit{oh}}
\def\Spacesstack{\mathcal{S}\!\mathit{paces}}
\def\Quotfunctor{\mathrm{Quot}}
\def\Hilbfunctor{\mathrm{Hilb}}
\def\Curvesstack{\mathcal{C}\!\mathit{urves}}
\def\Polarizedstack{\mathcal{P}\!\mathit{olarized}}
\def\Complexesstack{\mathcal{C}\!\mathit{omplexes}}
% \Pic is the operator that assigns to X its picard group, usage \Pic(X)
% \Picardstack_{X/B} denotes the Picard stack of X over B
% \Picardfunctor_{X/B} denotes the Picard functor of X over B
\def\Pic{\mathop{\mathrm{Pic}}\nolimits}
\def\Picardstack{\mathcal{P}\!\mathit{ic}}
\def\Picardfunctor{\mathrm{Pic}}
\def\Deformationcategory{\mathcal{D}\!\mathit{ef}}
\usepackage{cleveref}
\crefname{originalreference}{original reference}{original references}
\Crefname{originalreference}{Original reference}{Original references}
\makeatletter
\newlabel{definition-set-theoretically-invariant}{{048V}{source}{}{sourcepointer-1}{}}
\newlabel{definition-set-theoretically-invariant@cref}{{[originalreference][0][]048V}{[1][0][]source}{}{}{}}
\makeatother
\crefname{originallemma}{original lemma}{original lemmas}
\Crefname{originallemma}{Original lemma}{Original lemmas}
\makeatletter
\newlabel{schemes-lemma-immersion-when-closed}{{01IQ}{source}{}{sourcepointer-2}{}}
\newlabel{schemes-lemma-immersion-when-closed@cref}{{[originallemma][0][]01IQ}{[1][0][]source}{}{}{}}
\makeatother
\title{Set-theoretically invariant morphisms into locally separated targets are invariant when the relation is reduced}
\author{The Stacks Project Authors}
\date{stacks-project:groupoids-quotients.tex:local-snapshot-20261012; GNU FDL 1.2 or later}
\begin{document}
\maketitle
\noindent\textit{Editorial note.} All delivered mathematical text is the author's own, copied from the pinned Stacks Project chapter file groupoids-quotients.tex; no statement or proof step was rewritten, reordered or removed.\par\medskip
\noindent\textit{Editorial note.} The counted proof is the author's complete argument: The reducedness of R is essential; without it the bijective closed immersion Z to R need not be an isomorphism.\par\medskip
\noindent\textit{Editorial note.} The label lemma-set-theoretically-invariant-invariant-when-reduced names the counted statement, and the proof environment immediately following it in the source is delivered verbatim as its proof.\par\medskip
\noindent\textit{Editorial note.} The counted statement is carried by the author's own lemma environment in the chapter source.\par\medskip
\noindent\textit{Allowed standard prerequisites (editorial).}\begin{itemize}
\item Chapter Definition definition-set-theoretically-invariant (navigation stub): set-theoretically R-invariant means U(k) to X(k) equalizes s and t for every algebraically closed k over B; equivalently the equalizer Z below maps surjectively onto R on topological spaces, since every geometric point of R lifts to a k-point
\item X locally separated over B: the diagonal of X over B is an immersion, so the equalizer Z of the two maps from R to X times over B X, taken over the diagonal, is an immersion into R
\item Schemes, Lemma schemes-lemma-immersion-when-closed: an immersion whose image is closed is a closed immersion; in particular a bijective immersion is a closed immersion
\item A bijective closed immersion into a reduced algebraic space or scheme is an isomorphism: its defining quasi-coherent ideal sheaf consists of nilpotent sections and hence vanishes on a reduced target
\end{itemize}\medskip
\ifdefined\hypertarget\hypertarget{sourcepointer-1}{}\fi
\noindent\textit{Source reference.} Original reference 048V: sources/stacks-groupoids-quotients/pinned-original-source.tex, source lines 953--953. Prerequisite pointer invoked at 1002; it names a standard fact of the same chapter that the counted argument uses, and every step of the counted proof built on it is retained verbatim in the delivered body.\par\medskip
\ifdefined\hypertarget\hypertarget{sourcepointer-2}{}\fi
\noindent\textit{Source reference.} Original lemma 01IQ: sources/stacks-chow/tags.txt, source lines 1987--1987. Prerequisite pointer invoked at 1021; it names a standard fact of a sibling chapter that the counted argument uses, and every step of the counted proof built on it is retained verbatim in the delivered body.\par\medskip
\begin{lemma}
\label{lemma-set-theoretically-invariant-invariant-when-reduced}
In the situation of Definition \ref{definition-set-theoretically-invariant}.
Let $\phi : U \to X$ be a morphism of algebraic spaces over $B$.
Assume
\begin{enumerate}
\item $\phi$ is set-theoretically $R$-invariant,
\item $R$ is reduced, and
\item $X$ is locally separated over $B$.
\end{enumerate}
Then $\phi$ is $R$-invariant.
\end{lemma}

\begin{proof}
Consider the equalizer
$$
Z = R \times_{(\phi, \phi) \circ j, X \times_B X, \Delta_{X/B}} X
$$
algebraic space. Then $Z \to R$ is an immersion by assumption (3).
By assumption (1) $|Z| \to |R|$ is surjective. This implies that
$Z \to R$ is a bijective closed immersion (use
Schemes, Lemma \ref{schemes-lemma-immersion-when-closed})
and by assumption (2) we conclude that $Z = R$.
\end{proof}

\end{document}
